%!TEX encoding = UTF8
%!TEX root = 0-notes.tex

\chapter*{Préface}


\section*{Syllabus}

\subsection*{Description}

L’objectif de cette UE est double : une remise à niveau “sortie lycée” pour des étudiantes
et étudiants n’ayant pas suivi la spécialité mathématiques en terminale (mais pouvant avoir
suivi l’option mathématiques complémentaires) ainsi que de faire retravailler certains concepts
d’analyse du lycée, en les approfondissant, et en développant la pratique du calcul et 
l’interprétation des calculs. L’accent est mis sur la pratique et l’interprétation des calculs et non
sur la maîtrise des concepts théoriques (qui seront plus largement abordés dans d’autres UE).

\subsection*{Objectifs}

L’objectif est de faire travailler :
\begin{itemize}
	\item Traduction en équations et résolution d’équations.
	\item Études de fonctions.
	\item Rédaction et organisation des calculs, résolutions d’équations, études de fonctions.
	\item Interprétation géométrique des résultats, tracé de courbes.
	\item Articulation raisonnement et calcul.
\end{itemize}

Le programme donne un cadre pour atteindre cet objectif, tout en introduisant les 
contenus nécessaires. Il ne s’agit de travailler ces contenus avec des définitions “type” lycée, 
puisqu’ils seront repris, de façon approfondie dans l’UE Analyse I notamment.
Ces contenus sont travaillés et approfondis à partir de deux thèmes : étude élémentaires
de courbes et de surfaces, et équations différentielles.
Il ne s’agit pas de développer des outils théoriques ou des techniques avancées, mais
d’étudier et comprendre les objets à travers l’étude d’exemples.

\subsection*{Pré-requis nécessaires}

Programme de mathématiques du lycée, a minima spécialité de première et de préférence
option mathématiques complémentaires.

\subsection*{Pré-requis recommandés}

Non pertinent (UE de remise à niveau).

%\newcounter{preface}
\begin{Exercise}[label=1]%[counter=preface]
	Exemple d'exercice d'application directe.
\end{Exercise}
\begin{Exercise}[difficulty=1, label=2]%, counter=preface]
	Exemple de problème.
\end{Exercise}
\begin{Exercise}[difficulty=2, label=3]%, counter=preface]
	Exemple de problème d'approfondissement.
\end{Exercise}
\shipoutExercise
\begin{Answer}[ref=1]
	Solution 1.
\end{Answer}
\begin{Answer}[ref=2]
	Solution 2.
\end{Answer}
\begin{Answer}[ref=3]
	Solution 3.
\end{Answer}
\begin{multicols}{3}
\shipoutAnswer
\end{multicols}
